\begin{table}[!ht]
\centering
\caption{Two published expressions against the property the same paragraph claims for them. \textbf{Above}: the contrast term of the Qaroush visual-attention objective, over a sweep in which contrast rises by construction (25 levels, seed 42). Histogram equalisation drives every frame towards the same output standard deviation, so the published difference measures distance from maximum contrast: it runs backwards, and so does the signed form. Only the ratio, which the equation does not write, agrees. \textbf{Below}: the brightness term of the Korolkov module, over a sweep in which exposure rises by construction (21 levels). The source states that it ``penalizes under- and over-exposed frames'', which requires a non-monotone function; the mean luminance it specifies is monotone and enters the module with weight $+0.3$, so a brighter frame always scores higher. \textbf{No corpus is involved: both inversions are properties of the expressions.}}
\label{tab:ct}
\footnotesize
\setlength{\tabcolsep}{3.5pt}
\renewcommand{\arraystretch}{1.35}
\begin{tabular}{@{}>{\raggedright\arraybackslash}p{4.0cm} l r@{}}
\toprule
\rowcolor{headerblue}
\hd{Reading} & \hd{Expression} & \hd{$\rho$ with the property} \\
\midrule
\multicolumn{3}{@{}l}{\textit{Qaroush contrast, against contrast (Eq.~8)}} \\
\quad Eq.~8 as published & $\frac{1}{K}\sum_k\bigl|\sigma_p(\mathrm{HE}(f_k))-\sigma_p(f_k)\bigr|$ & \textcolor{accentred}{\textbf{$-1.0000$}} \\
\rowcolor{rowlight}\quad without the absolute value & $\frac{1}{K}\sum_k\bigl(\sigma_p(\mathrm{HE}(f_k))-\sigma_p(f_k)\bigr)$ & \textcolor{accentred}{\textbf{$-1.0000$}} \\
\quad ratio, which the source does not write & $\frac{1}{K}\sum_k \sigma_p(f_k)/\sigma_p(\mathrm{HE}(f_k))$ & $+1.0000$ \\
\midrule
\multicolumn{3}{@{}l}{\textit{Korolkov brightness, against exposure}} \\
\quad as published & $\frac{1}{|I|}\sum_{x,y} L(I)_{x,y}$ (LAB) & $+1.0000$ \\
\rowcolor{rowlight}\quad what ``penalizes'' would need & $-\bigl|L-127.5\bigr|$ & \textcolor{accentred}{\textbf{$-.4883$}} \\
\bottomrule
\end{tabular}
\end{table}
